% =============================================================================
%  ECMLPKDD_DatasetExps.tex
%
%  Holding file for experimental content that was produced on a previous
%  dataset release (the ECMLPKDD-era TurboSens1 and an interim
%  scenario-3 generation that is no longer part of the released benchmark).
%
%  These paragraphs were lifted out of the main NeurIPS 2026 submission on
%  2026-04-26 to keep the paper focused on TurboSens3 and TurboSens4 only.
%  Numbers, component labels, and figures here refer to data that no longer
%  ships with the repository, so they should not be re-introduced into the
%  main body without re-running the corresponding probes on the current
%  TurboSens releases.
%
%  Sections preserved verbatim:
%    A.  Synthetic OOD scenario descriptions (formerly main-body §Task 4)
%    B.  §Results subsection (Tasks 1, 2, 2bis, 3, 4 prose)
%    C.  Old Task 1 HI estimation table (tab:task1)
%    D.  Old Task 3 forecasting table (tab:task3)
%    E.  Older paper_results_section.tex draft with the (w, H) probe sweep
% =============================================================================

\section*{A. Synthetic OOD scenario descriptions (legacy)}

The three paragraphs below described the synthetic OOD construction that the
older experimental pipeline used. The current paper instead evaluates OOD
generalisation on the held-out \texttt{test\_hard} split of TurboSens4
(different episode seeds, drawn from the same region prior with a higher
extreme-degradation rate). The synthetic constructions below are kept for
reference only.

\paragraph{OOD scenario 1: premature maintenance.}
A maintenance action (HI restoration) is triggered while the engine health
is still near nominal, a configuration never present in the training set
where maintenance always follows significant degradation.
The predictor, calibrated to slow monotonic drift, expects continuation of
near-nominal dynamics; the sudden restoration jump produces a spike in
$\mathcal{S}$ localised to the event timestep.

\paragraph{OOD scenario 2: accelerated degradation.}
Degradation-rate parameters are set to values outside the training range
(multiplied by factors of $3\times$ and $5\times$).
The latent trajectory traverses the learned degradation manifold at an
anomalous speed; the predictor underestimates the per-step displacement,
causing $\mathcal{S}$ to grow proportionally to the rate multiplier.

\paragraph{OOD scenario 3: correlated component failure.}
Two turbofan components whose degradation is statistically independent in
the training distribution (fan and turbine) are forced to degrade
simultaneously at elevated rate.
Even if each component's individual rate lies within the training range,
the joint dynamics occupy a low-density region of the latent space,
detectable via the Mahalanobis distance
$d_M(\mathbf{z}_t) = \sqrt{(\mathbf{z}_t - \boldsymbol{\mu})^\top
                     \boldsymbol{\Sigma}^{-1}
                     (\mathbf{z}_t - \boldsymbol{\mu})}$
computed from the training embedding distribution.


% -----------------------------------------------------------------------------
\section*{B. Legacy main-body §Results subsection}

The four paragraphs below report HI-state estimation, velocity, forecasting
and OOD detection on a JEPA model trained on a previous scenario release
that is no longer in the repository. The mean $R^2$ values
($\bar R^2 = 0.83$ at $w{=}1$, $0.85$ at $w{=}10$) and the $\delta(\tau)$
peak around $\tau \approx 6$ are from that older snapshot. They predate
the dragon / OT campaign on TurboSens4 reported in the current main body.

\paragraph{Task 1: Health-Indicator State Estimation.}
Table~\ref{tab:legacy_task1} reports per-component $R^2$, RMSE, and Pearson-$r$ for the HI state estimation probe trained on frozen embeddings.
The snapshot encoder ($w{=}1$) achieves a mean $R^2$ of $0.834$ across all ten turbofan health-indicator dimensions, with the single-embedding probe already recovering high-amplitude components nearly perfectly ($R^2 = 0.956$ for \texttt{H/Eff}, $R^2 = 0.935$ for \texttt{Fan/Wc}).
The temporally aggregated encoder ($w{=}10$) improves the mean to $0.852$, with gains concentrated in low-amplitude, slowly varying components: \texttt{Bst/Wc} improves by $+0.10$ and \texttt{TrbH/Eff} by $+0.06$, while high-amplitude components remain essentially unchanged.
This pattern is consistent with the TemporalAggregator acting as a noise-smoothing mechanism: components whose instantaneous sensor signal is weak benefit from averaging over a short trajectory window, whereas already-discriminative components saturate regardless. The lowest-performing dimension, \texttt{TrbL/Wc} ($R^2 \approx 0.71$), corresponds to a slow, low-amplitude degradation mode that is hard to separate from sensor noise at any window size.

\paragraph{Task 2: Degradation Velocity Estimation.}
The velocity probe (predicting $\Delta\mathbf{s}_t = \mathbf{s}_t - \mathbf{s}_{t-1}$ from a single embedding) yields $R^2 \approx -0.02$ for both model variants, with the probe converging to predicting the mean (near zero) within 21 epochs. This result is expected and informative.
The per-step degradation increment is approximately $2.3 \times 10^{-5}$ in normalised units, five orders of magnitude smaller than the HI amplitude, with variance so low that the RMSE of a constant-mean predictor equals the empirical standard deviation.
The JEPA objective only requires predicting the \emph{position} of the next latent state, not the velocity of transition; accordingly, the encoder has no gradient pressure to encode degradation rate in its geometry. We evaluate the same probe with trajectory windows of $\mathrm{seq\_len} \in \{1, 10, 50\}$:
a probe attending over a window of consecutive embeddings can compute a finite-difference estimate of the local slope, which is the natural
information source for velocity estimation. Whether this suffices to push $R^2$ into positive territory is an open question we report in Table~\ref{tab:legacy_task3}.
The maintenance alarm (Task~2b), scored by the predicted $\|\Delta\hat{\mathbf{s}}_t\|_2$, yields AUC$=0.50$ at all horizons when the velocity probe is constant, confirming that snapshot embeddings carry no direct alarm signal.

\paragraph{Task 2 (bis): Degradation Archetype Classification.}
Simply from some sequence, predict the current degradation archetype.

\paragraph{Task 3: Latent Forecasting and Action Divergence.}
Figure~\ref{fig:legacy_task3} reports the mean HI RMSE of the zero-action autoregressive rollout as a function of forecast horizon $\tau$. We distinguish \emph{clean} trajectories (no maintenance event in $(t, t{+}\tau]$) from \emph{event} trajectories (at least one maintenance event), and define the action-divergence gap $\delta(\tau) = \mathrm{RMSE}_\mathrm{event}(\tau) - \mathrm{RMSE}_\mathrm{clean}(\tau)$.

\begin{figure}[h]
    \centering
    \includegraphics[width=0.9\linewidth]{figures/fig_task3_curves.pdf}
    \caption{Latent forecasting comparison of window sizes for le-wm over latent forecasting (Task 3). Legacy figure on a previous scenario release.}
    \label{fig:legacy_task3}
\end{figure}

Three findings stand out. First, the one-step clean RMSE of the world model matches the direct encoder RMSE from Task~1 to within numerical precision: $0.00199$ vs.\ $0.00198$ for $w{=}1$, and $0.00186$ vs.\ $0.00187$ for $w{=}10$.
The predictor adds negligible error at $\tau{=}1$ for smooth degradation trajectories; the encoder representation is the performance bottleneck, not the dynamics model.

Second, the $w{=}10$ model consistently outperforms $w{=}1$ on clean trajectories at all horizons, with RMSE at $\tau{=}50$ of $0.00480$ versus $0.00572$ (a $16\%$ reduction). Temporal aggregation at the encoder level propagates useful trajectory direction information into the initial latent state, giving the predictor a better-informed context for multi-step rollout.

Third, the action-divergence gap $\delta(\tau)$ peaks at $\tau \approx 6$ (gap $\approx +0.0031$ for $w{=}1$, $+0.0029$ for $w{=}10$) and decays monotonically toward zero by $\tau{=}50$. The peak reflects the average proximity of a sampled starting frame to the next maintenance event in the dataset. The decay is a statistical dilution effect: at large $\tau$, the event set is dominated by trajectories whose maintenance event falls late in the window, contributing minimal extra error at the evaluated prediction target. The peak gap ($\approx 0.003$) is larger than the per-step degradation drift at any horizon, establishing maintenance transitions as the dominant source of forecast error. This motivates action-conditioned world modelling as a direct and quantified next step: closing the gap would require the predictor to condition on the binary maintenance signal, which is precisely the information discarded in the zero-action baseline.

We evaluate three unsupervised OOD detectors on frozen JEPA representations. Distribution-based methods (Mahalanobis, k-NN) achieve moderate AUC (0.69-0.77), demonstrating that JEPA embeddings carry distributional information about the training regime. Notably, prediction error (Surprise) is an unreliable OOD signal — OOD trajectories are often more predictable than training ones because fast monotonic degradation is locally smooth. This motivates future work on representations explicitly trained with anomaly-sensitivity objectives.


% -----------------------------------------------------------------------------
\section*{C. Legacy Task 1 table}

%% ── Table 1: Task 1 HI State Estimation ─────────────────────────────────────
\begin{table}[h]
\centering
\caption{%
  Legacy Task~1: Health-indicator state estimation.
  A TransformerProbe is trained on frozen embeddings to predict
  each of the ten turbofan HI dimensions.
  Results are shown for the snapshot encoder ($w{=}1$) and the
  temporally aggregated encoder ($w{=}10$). Numbers are from a
  previous scenario release.
}
\label{tab:legacy_task1}
\setlength{\tabcolsep}{5pt}
\begin{tabular}{lcccccc}
\toprule
& \multicolumn{3}{c}{$w = 1$} & \multicolumn{3}{c}{$w = 10$} \\
\cmidrule(lr){2-4} \cmidrule(lr){5-7}
Component & $R^2$ & RMSE & Pearson & $R^2$ & RMSE & Pearson \\
\midrule
Bst/Eff    & 0.847 & 0.00185 & 0.926 & 0.850 & 0.00183 & 0.930 \\
Bst/Wc     & 0.762 & 0.00233 & 0.883 & 0.865 & 0.00176 & 0.939 \\
Fan/Eff    & 0.782 & 0.00217 & 0.888 & 0.785 & 0.00216 & 0.891 \\
Fan/Wc     & 0.935 & 0.00107 & 0.969 & 0.954 & 0.00090 & 0.979 \\
H/Eff      & 0.956 & 0.00162 & 0.978 & 0.957 & 0.00159 & 0.978 \\
H/Wc       & 0.863 & 0.00285 & 0.930 & 0.874 & 0.00274 & 0.935 \\
TrbH/Eff   & 0.792 & 0.00202 & 0.890 & 0.850 & 0.00172 & 0.924 \\
TrbH/Wc    & 0.869 & 0.00145 & 0.933 & 0.855 & 0.00152 & 0.926 \\
TrbL/Eff   & 0.819 & 0.00195 & 0.906 & 0.818 & 0.00196 & 0.906 \\
TrbL/Wc    & 0.713 & 0.00249 & 0.845 & 0.714 & 0.00248 & 0.846 \\
\midrule
\textbf{Mean} & \textbf{0.834} & \textbf{0.00198} & \textbf{0.915}
              & \textbf{0.852} & \textbf{0.00187} & \textbf{0.925} \\
\bottomrule
\end{tabular}
\end{table}


% -----------------------------------------------------------------------------
\section*{D. Legacy Task 3 table}

%% ── Table 2: Task 3 Summary ──────────────────────────────────────────────────
\begin{table}[h]
\centering
\caption{%
  Legacy Task~3: Latent forecasting RMSE at selected horizons $\tau$.
  \emph{Clean}: no maintenance event in $(t, t{+}\tau]$.
  \emph{Event}: at least one maintenance event in $(t, t{+}\tau]$.
  \emph{Gap}: $\delta(\tau) = \mathrm{RMSE}_\mathrm{event} -
  \mathrm{RMSE}_\mathrm{clean}$, quantifying the action-divergence cost of the zero-action forecaster. All values are mean HI RMSE $\times 10^{3}$.
  Numbers are from a previous scenario release.
}
\label{tab:legacy_task3}
\setlength{\tabcolsep}{5pt}
\begin{tabular}{r ccc ccc}
\toprule
& \multicolumn{3}{c}{$w = 1$} & \multicolumn{3}{c}{$w = 10$} \\
\cmidrule(lr){2-4} \cmidrule(lr){5-7}
$\tau$ & Clean & Event & Gap & Clean & Event & Gap \\
\midrule
 1  & 1.99 & 4.52 & $+$2.54 & 1.86 & 4.28 & $+$2.42 \\
 5  & 2.07 & 5.10 & $+$3.04 & 1.91 & 4.75 & $+$2.83 \\
10  & 2.22 & 5.04 & $+$2.83 & 2.03 & 4.64 & $+$2.61 \\
20  & 3.04 & 5.03 & $+$1.99 & 2.57 & 4.43 & $+$1.86 \\
30  & 4.21 & 5.26 & $+$1.05 & 3.47 & 4.49 & $+$1.01 \\
50  & 5.72 & 5.89 & $+$0.17 & 4.80 & 4.75 & $-$0.05 \\
\bottomrule
\end{tabular}
\end{table}


% -----------------------------------------------------------------------------
\section*{E. Older paper\_results\_section.tex draft (full $w \times H$ sweep)}

The block below was previously held in \texttt{paper\_results\_section.tex}
at the repo root. It reports a six-variant probe sweep across two
architectures (JEPA, AR-LSTM) at three encoder window sizes
$w \in \{1, 50\}$ and three predictor history lengths
$H \in \{1, 50, 100\}$. The numbers refer to the same legacy scenario
release as section~B above.

\begin{table}[h]
\centering
\caption{%
  \textbf{Legacy Task 1 — HI State Estimation, $w \times H$ sweep.}
  Frozen encoder probe trained on embeddings from six model variants.
  Means are macro-averaged over the ten HI dimensions.
  Best result per metric in \textbf{bold}.
}
\label{tab:legacy_task1_sweep}
\setlength{\tabcolsep}{5pt}
\begin{tabular}{llcc rrr cc}
\toprule
\multirow{2}{*}{Architecture} & \multirow{2}{*}{Config} &
  \multirow{2}{*}{$w$} & \multirow{2}{*}{$H$} &
  \multicolumn{3}{c}{Mean (all 10 dims)} &
  \multicolumn{2}{c}{Hardest dims} \\
\cmidrule(lr){5-7}\cmidrule(lr){8-9}
 & & & & R$^2$ & RMSE\textsuperscript{\dag} & Pearson-$r$ & TrbH/Eff & TrbL/Wc \\
\midrule
AR-LSTM & $w{=}1,H{=}1$  & 1 &  1 & \textbf{0.792} & 2.23 & \textbf{0.892} & 0.782 & 0.687 \\
AR-LSTM & $w{=}1,H{=}50$ & 1 & 50 & 0.776          & 2.31 & 0.887          & 0.762 & 0.661 \\
\midrule
JEPA    & $w{=}1,H{=}50$ & 1 &  50 & 0.763          & 2.42 & 0.878          & 0.760 & 0.683 \\
JEPA    & $w{=}1,H{=}100$& 1 & 100 & 0.785          & 2.26 & 0.891          & 0.769 & 0.691 \\
JEPA    & $w{=}50,H{=}1$ & 50&   1 & 0.709          & 2.13 & 0.857          & 0.643 & 0.532 \\
JEPA    & $w{=}50,H{=}50$& 50&  50 & 0.736          & \textbf{2.03} & 0.874 & 0.634 & 0.622 \\
\bottomrule
\end{tabular}\\
\textsuperscript{\dag}RMSE in units of normalised HI ($\times 10^{-3}$).
\end{table}

% -----------------------------------------------------------------------------
% End of file. Anything below this line is for the maintainer's notes.
% -----------------------------------------------------------------------------
